\chapter{The accessible genome at Single-Cell Resolution}
The work performed in this thesis aims to advance a technical aspect (single-cell sequencing) of an important part (chromatin accessibility) of the gene regulatory puzzle. This introductory chapter will attempt to build the foundation for the explorations that follow by providing a historical overview of chromatin accessibility and frame it within the ongoing research on decoding gene regulation. Then, we frame the state of the art of single-cell chromatin sequencing technologies when I began this work, and delve into the capabilities and limitations of several breakthrough techniques. This context serves to highlight the importance of the subject matter, and to motivate our subsequent discussions and findings.

\section{A brief historical overview of cellular genetics}
\begin{wrapfigure}[9]{R}{0.15\textwidth} %this figure will be at the right
	\centering
	\raisebox{0pt}[\dimexpr\height-1\baselineskip\relax]{\includegraphics[width=0.15\textwidth]{theodor_schwann.jpg}}%
	\caption*{T. Schwann}
\end{wrapfigure}
The word "cell" was first used by Robert Hooke to describe the hexagonal cell walls in a thin slice of cork in the 1600s\sjcitep[,]{hookeMicrographia1665} but it took several more centuries of research until the concept of cells as the foundational building block of life began to take shape. Only in the mid 1800s, Theodor Schwann showed that animals, just like plants, were composed of cells\sjcitep[,]{schwannMikroskopischeUntersuchungenUber1839} and Rudolf Virchow and Robert Remak later added that cells were the product of division of pre-existing cells\sjcitep[.]{rudolfCellularpathologieIhrerBegrundung1858} While studying this process of division in 1879, Walther Flemming first used the term "chromatin" to describe the filament-like threads between two dividing cells in 1879\sjcitep[.]{flemmingBeitrageZurKenntniss1879} Flemming's observations did not differ much from his colleague's Wilhelm Waldeyer's description of "chromosomes"\sjcitep[.]{deichmannChromatinItsHistory2015} At the time, the physical transfer of genetic information was not yet widely accepted, and the exact function of these chromosomes, and how genetic information was stored, was still unknown. Two breakthrough discoveries were made in the lab of Thomas Hunt Morgan, who in 1908 began studying inheritance by irradiating \textit{Drosophila melanogaster} (fruit fly) and cross-breeding the resulting mutants.

\begin{wrapfigure}[9]{L}{0.15\textwidth} %this figure will be at the right
	\centering
	\includegraphics[width=0.15\textwidth]{thomas_morgan.jpg}
	\caption*{T. Morgan}
\end{wrapfigure}

After several years without results, Morgan observed that some mutant traits, such as the flies' eye colours, could be predictably inherited\sjcitep[.]{morganRandomSegregationCoupling1911} More than half a century earlier, Moravian friar Gregor Mendel had demonstrated the exact same principles using \textit{Pisum sativum} (the common sweet pea), but his work was lost and disputed upon rediscovery. Using his new findings, Morgan verified Mendel's work and formally defined the theory of Mendelian inheritance in 1915\sjcitep[.]{morganMechanismMendelianHeredity1915} A second breakthrough came when Morgan demonstrated that the chromosomes first observed by Flemming and Waldeyer were in fact the physical bodies responsible for the transfer of phenotypic traits, confirming earlier ideas that chromosomes carried genes which encoded heritable traits. Together, Morgan's findings inaugurated the 20\textsuperscript{th} century's era of modern molecular genetics and set the stage for many of the major discoveries that ensued. However, until this point, a handful of fundamental questions about the nature of cells and life in general had kept the minds of biologists occupied, but remained completely unanswered:

\begin{enumerate}[noitemsep]
	\item How can a single fertilized egg cell self-organise into an array of thousands of different cell types?
	\item Where is the information that steers this process encoded, and how is it interpreted?
	\item How can this process go wrong, leading to disease?
\end{enumerate}

Several groups were working on these issues in the early 20\textsuperscript{th} century. It had been known that a small portion of the vertebrate embryo acted as an "organizing center", seemingly directing nearby tissues throughout the early development of the embryo. Hermann Spemann together with his student Hilde Mangold discovered that, when transplanted from one \textit{Triturus} (newt) embryo to another, this "organizer" could induce the formation of a secondary embryo opposite to the other\sjcitep[.]{spemannUberInduktionEmbryonenanlagen1924} Multiple labs raced to purify and characterize the compound responsible for these processes in an "organizer gold rush"\sjcitep[,]{slackConradHalWaddington2002} but without success. We now know that the responsible compounds are growth factors, which are extremely potent and already active at nanogram quantities - far below the detection limit of that era's technology.

\begin{wrapfigure}[9]{L}{0.15\textwidth} %this figure will be at the right
	\centering
	\raisebox{0pt}[\dimexpr\height-1\baselineskip\relax]{{}\includegraphics[width=0.15\textwidth]{conrad_waddington.png}}
	\caption*{C. Waddington}
\end{wrapfigure}

At the time, the link between embryonic development and the new insights in genetics was not yet understood, even dismissed. This changed when a 25 year old Conrad Waddington, who had been performing grafting experiments on organizers in mammals and birds, visited Morgan's fly lab in 1935, where he intensely studied \textit{Drosophila}'s embryonic wing development. Waddington discovered that specific, heritable mutations could disrupt wing development at different stages, or even lead to the growth of terminal leg segments in place of the fly's antenna. Based on these findings, he wrote his seminal work \textit{Organisers and Genes}, which introduced the notion that embryonic development was a decision tree under the control of genes, and that these decisions together formed an "epigenetic landscape"\sjcitep[.]{waddingtonOrganisersGenes1940} Around that time, the role of DNA as a store of genetic information was being established\sjcitep[,]{averyStudiesChemicalNature1944} and the exact molecular mechanism was soon resolved by Watson, Crick and Franklin\sjcitep[.]{watsonMolecularStructureNucleic1953,watsonGeneticalImplicationsStructure1953}. Waddington extensively elaborated on his concept of the epigenetic landscape in his second book, \textit{The Strategy of Genes}\sjcitep[,]{conradStrategyGenes1957} in 1957. Here, he formulated the landscape as a multidimensional metabolic state space. A ball, representing a cell, rolls down this landscape during development, passing several bifurcations which ultimately determine its "fate" (\autoref{fig:waddington_landscape}a). Beneath the landscape, a series of tension ropes, symbolizing genes, forcefully pull the surface into shape (\autoref{fig:waddington_landscape}b).

\begin{figure}[h]
    \centering
    \includegraphics[width=0.75\textwidth]{waddington_landscape.png}
	\captionsetup{width=0.75\linewidth}
    \caption{Waddington's epigenetic landscape. }
    \label{fig:waddington_landscape}
\end{figure}

A great milestone that played exactly into Waddington's ideas was the discovery of the \textit{lac} operon. Jacques Monod, \textit{another} student of Thomas Morgan, had observed that bacteria grown in a medium containing both glucose and lactose would preferentially feed on glucose, and only begin feeding on lactose once glucose had run out. Further experiments together with François Jacob led to the understanding that bacteria only began producing galactose, the enzyme necessary to metabolise lactose, both in the presence of lactose and in the absence of glucose. The pair discovered that the production of this enzyme was under the control of a repressor, also a protein, which became inactive in the presence of galactose. Together, they published this new theory in 1961\sjcitep[,]{jacobGeneticRegulatoryMechanisms1961} in a review which features the first diagram of a gene regulatory network (\autoref{fig:lac_repressor}).

\begin{figure}[h]
    \centering
    \includegraphics[width=0.75\textwidth]{lac_repressor.png}
	\captionsetup{width=0.75\linewidth}
    \caption{Jacob \& Monod's proposed model for the regulation of protein synthesis. A regulator gene produces a repressor, which represses expression of genes A and B unless a specific metabolite is present. Reproduced from \notefullcite{jacobGeneticRegulatoryMechanisms1961}.}
    \label{fig:lac_repressor}
\end{figure}

Earlier that year, Jacob, together with Sidney Brenner, had published the first observations of an unstable "messenger" RNA\sjcitep[.]{brennerUnstableIntermediateCarrying1961} This mRNA was clearly depicted in Jacob \& Monod's gene regulatory network as the intermediate between genes and proteins. Moreover, the diagram explicitly showed the repressor itself as a protein originating from a gene. Genes thus controlled other genes - exactly how Waddington had imagined. At this point, several pieces of the puzzle began to fall into place. It became clear that the developmental questions we posed earlier could all be reduced to the concept and foundational nature of \textit{gene regulation}.

We now understand that fine-tuned gene regulation explains how one cell can give rise to many, that this program is embedded in our genomes through a "regulatory code", and that aberrant gene regulation can explain many of our age's most important diseases. Much progress has been made in unraveling the mechanics of gene regulation. We have models of gene transcription and translation, cell signaling, the roles of several classes of proteins, and so on. However, full and deterministic models which can exhaustively answer our core questions remain elusive, almost unimaginable.

\section{Next generation sequencing techniques}
A second wave of relief: NGS
\begin{itemize}
	\item microarrays
	\item bulk RNA-seq
	\item FAIRE-seq, DNAse-seq etc
	\item bulk ATAC-seq
\end{itemize}

\section{The central dogma of single-cell sequencing}
Cell type diversity underlies \\
Single-cell sequencing was \supershortnotecite{noauthorslistedMethodYear20132014}

\begin{itemize}
	\item NGS sequencing by itself provided a lot of relief, but not quite
\end{itemize}

\lipsum[1]

\section{Generation 0 - The early years}
The idea of extracting material from single cells is not new.
\begin{itemize}
	\item Not a new idea
	\item Eberwine
\end{itemize}

By isolating, sorting and sequencing the genome of 52 cells from a primary breast tumour and 48 cells from its associated liver metastasis, Navin et al. showed that a single cell from the primary tumor mass seeded the metastatic tumour\sjcitep[.]{navinTumourEvolutionInferred2011}

\lipsum[1]

\section{Generation 1 - Small steps}
\begin{itemize}
	\item Method of the year 2013
	\item First generation technologies such as CEL-seq, SMART-seq, inDrop, Drop-seq
\end{itemize}

\lipsum[1]

\section{Generation 2 - Learning how to run}
\begin{itemize}
	\item Wide spread of 10x
	\item Addition of ATAC-seq and possibly others
\end{itemize}

\lipsum[1]

\section{Generation 3 - Large leaps}
\begin{itemize}
	\item Super large scale sequencing projects
	\item Multiomics techniques
	\item Method of the year 2019
\end{itemize}

\lipsum[1]

\section{Off into space}
\begin{itemize}
	\item Similar evolutions as SC techs themselves
	\item Method of the year
\end{itemize}

\lipsum[1]

\section{Computational models}
\begin{itemize}
	\item how the sequencing works
	\item count matrices in both rna and ATAC
	\item multiome added value examples
	\item using genomic information to assign samples
	\item atac: extra challenge
	\item AI models
\end{itemize}

%% This part makes a figure
\begin{figure*}[h]
	\dummyfig{Dummy Figure Label}
	\caption{Dummy figure caption}
	\label{fig:dummy}
\end{figure*}